\documentclass{article}
\usepackage{geometry}
\usepackage{enumitem}
\usepackage{hyperref}



\begin{document}

\section*{Peer Review Summary}
\begin{enumerate}
    \item \textbf{Coherent and Engaging Historical Narrative}
    \begin{enumerate}
        \item \textbf{Strengths:}
        \begin{enumerate}
            \item The cross-cultural comparisons (Chinese vs. Greek vs. Babylonian approaches) add depth and highlight the uniqueness of Chinese mathematical philosophy.
            \item The applied case studies (e.g., Dujiangyan’s hydraulic engineering, Han beacon towers) anchor abstract concepts in tangible historical practices.
        \end{enumerate}
        
        \item \textbf{Areas for Improvement:}
        \begin{enumerate}
            \item Smoother transitions are needed between sections (e.g., from Three Kingdoms proofs to Song-Yuan mathematics on Page 7).
            \item Expand on why Chinese mathematics prioritized practicality (e.g., bureaucratic demands for taxation, infrastructure projects).
        \end{enumerate}
    \end{enumerate}
    
    \item \textbf{Readability and Clarity}
    \begin{enumerate}
        \item \textbf{Strengths:}
        \begin{enumerate}
            \item The text is vivid and accessible, with phrases like “toolkit for building empires” enhancing engagement.
            \item Equations are formatted clearly and explained in context (e.g., hydraulic engineering formula on Page 8).
        \end{enumerate}
        
        \item \textbf{Areas for Improvement:}
        \begin{enumerate}
            \item Add visual aids for Liu Hui’s proof (Pages 4–5) to improve clarity.
            \item Define variables like \(k\) in the military defense formula \(c = \sqrt{a^2 + (k \cdot b)^2}\).
            \item Add diagrams for Liu Hui’s proof (e.g., labeled sketches of geometric transformations).
        \end{enumerate}
    \end{enumerate}
    
    \item \textbf{Integration of Mathematics and History}
    \begin{enumerate}
        \item \textbf{Strengths:}
        \begin{enumerate}
            \item Mathematics is well-attributed (e.g., Liu Hui’s "Chu-Ru Xiangbu" proof tied to Three Kingdoms-era priorities).
            \item Discussion of transmission routes (e.g., Tang Dynasty exchanges) links ideas to historical movements.
        \end{enumerate}
        
        \item \textbf{Areas for Improvement:}
        \begin{enumerate}
            \item Integrate primary sources like “Jayan Han slips” mentioned in your reflection.
            \item Explain derivations (e.g., Zhang Heng’s $\theta = \arctan\left(\frac{a}{b}\right)$).
        \end{enumerate}
    \end{enumerate}
    
\end{enumerate}


\end{document}